\section{Conclusion}
\label{sec:main-conclusion}

We introduced proximal balancing, which estimates the average treatment effect from a multi-dimensional proxy without designating proxy roles, solving an inverse problem, or modeling the hidden confounder. When a held-out block does not affect the treatment and preserves the outcome-relevant part of the confounder, adjusting for a representation that balances it identifies the effect (Thm.~\ref{thm:exact-proxy-identification}). When the valid block is unknown, a plurality vote over held-out blocks identifies the effect (Thm.~\ref{thm:unknown-valid-block-identification}); under approximate balance, the bias is bounded by a stability-weighted residual treatment discrepancy (Thm.~\ref{thm:approximate-proxy-bias-bound}); and our algorithm, PROBE (Algo.~\ref{alg:probe-search}), has finite-sample guarantees for representation learning, honest AIPW estimation, and search (Thms.~\ref{cor:representation-causal}--\ref{thm:probe-search-error}). In synthetic data with vector and image proxies, PROBE had the smallest error of all competing methods other than the oracle, and on real-world data it remained competitive without role information, while methods given wrong roles degraded sharply.
